% ========================================================
% BÀI TẬP TOÁN 5 - KẾT NỐI TRI THỨC VỚI CUỘC SỐNG
% Bài 30: Ôn tập số thập phân
% Định dạng: BTTeX Part III - Trắc nghiệm điền số ngắn
% ========================================================

% Câu 1
\begin{ex}
  Trong số thập phân $245{,}78$, chữ số $8$ có giá trị bằng bao nhiêu phần trăm?
  \shortans{8}
  \loigiai{
    Chữ số 8 ở hàng phần trăm, có giá trị 8 phần trăm.
  }
\end{ex}

% Câu 2
\begin{ex}
  Số thập phân gồm $3$ chục, $5$ đơn vị và $6$ phần mười là $35{,}6$. Tổng các chữ số của số đó là bao nhiêu?
  \shortans{14}
  \loigiai{
    $3 + 5 + 6 = 14$.
  }
\end{ex}

% Câu 3
\begin{ex}
  Viết phân số $\frac{4}{5}$ thành số thập phân ta được $0{,}8$. Lấy phần nguyên cộng chữ số hàng phần mười được bao nhiêu?
  \shortans{8}
  \loigiai{
    $0 + 8 = 8$.
  }
\end{ex}

% Câu 4
\begin{ex}
  Tìm chữ số $x$ lớn nhất để $4{,}x2 < 4{,}72$.
  \shortans{6}
  \loigiai{
    $x < 7 \implies x = 6$.
  }
\end{ex}

% Câu 5
\begin{ex}
  Tìm chữ số $y$ bé nhất để $19{,}8y > 19{,}85$.
  \shortans{6}
  \loigiai{
    $y > 5 \implies y = 6$.
  }
\end{ex}

% Câu 6
\begin{ex}
  Làm tròn số $45{,}67$ đến hàng đơn vị ta được số nào?
  \shortans{46}
  \loigiai{
    Chữ số hàng phần mười là 6 nên làm tròn lên thành 46.
  }
\end{ex}

% Câu 7
\begin{ex}
  Làm tròn số $128{,}34$ đến hàng đơn vị ta được số nào?
  \shortans{128}
  \loigiai{
    Chữ số hàng phần mười là 3 nên làm tròn xuống thành 128.
  }
\end{ex}

% Câu 8
\begin{ex}
  Đổi $5\text{ m } 6\text{ dm}$ thành mét dưới dạng số thập phân là $5{,}6\text{ m}$. Chữ số hàng phần mười là bao nhiêu?
  \shortans{6}
  \loigiai{
    Chữ số là 6.
  }
\end{ex}

% Câu 9
\begin{ex}
  Đổi $3\text{ tấn } 25\text{ kg}$ thành tấn dưới dạng số thập phân là $3{,}025\text{ tấn}$. Chữ số hàng phần trăm là bao nhiêu?
  \shortans{2}
  \loigiai{
    Chữ số hàng phần trăm là 2.
  }
\end{ex}

% Câu 10
\begin{ex}
  Có bao nhiêu số tự nhiên $x$ thỏa mãn $4{,}5 < x < 9{,}5$?
  \shortans{5}
  \loigiai{
    $x \in \{5, 6, 7, 8, 9\}$. Có 5 số.
  }
\end{ex}

% Câu 11
\begin{ex}
  Viết số thập phân $0{,}75$ dưới dạng phân số tối giản $\frac{a}{b}$. Mẫu số $b$ là bao nhiêu?
  \shortans{4}
  \loigiai{
    $0{,}75 = \frac{3}{4}$. Mẫu số là 4.
  }
\end{ex}

% Câu 12
\begin{ex}
  Viết hỗn số $4\frac{3}{10}$ thành số thập phân ta được $4{,}3$. Tích của phần nguyên và phần thập phân là bao nhiêu?
  \shortans{12}
  \loigiai{
    $4 \times 3 = 12$.
  }
\end{ex}

% Câu 13
\begin{ex}
  Tìm số tự nhiên bé nhất lớn hơn $19{,}99$. Số đó là bao nhiêu?
  \shortans{20}
  \loigiai{
    Số tự nhiên bé nhất lớn hơn $19{,}99$ là 20.
  }
\end{ex}

% Câu 14
\begin{ex}
  Tìm số tự nhiên lớn nhất bé hơn $35{,}01$. Số đó là bao nhiêu?
  \shortans{35}
  \loigiai{
    Số tự nhiên lớn nhất bé hơn $35{,}01$ là 35.
  }
\end{ex}

% Câu 15
\begin{ex}
  Số thập phân $15{,}040$ sau khi bỏ các chữ số 0 vô nghĩa ở tận cùng bên phải phần thập phân có bao nhiêu chữ số ở phần thập phân?
  \shortans{2}
  \loigiai{
    $15{,}04$ có 2 chữ số ở phần thập phân.
  }
\end{ex}

% Câu 16
\begin{ex}
  Đổi $0{,}8\text{ m}^2$ thành bao nhiêu đề-xi-mét vuông?
  \shortans{80}
  \loigiai{
    $0{,}8 \times 100 = 80\text{ dm}^2$.
  }
\end{ex}

% Câu 17
\begin{ex}
  Đổi $2{,}5\text{ ha}$ thành bao nhiêu mét vuông? (Điền số nghìn mét vuông).
  \shortans{25}
  \loigiai{
    $2{,}5 \times 10\,000 = 25\,000\text{ m}^2 = 25$ nghìn mét vuông.
  }
\end{ex}

% Câu 18
\begin{ex}
  Tìm chữ số $a$ để $0{,}a = \frac{7}{10}$.
  \shortans{7}
  \loigiai{
    $a = 7$.
  }
\end{ex}

% Câu 19
\begin{ex}
  Từ các chữ số $2; 4; 6$, viết số thập phân lớn nhất có phần nguyên là một chữ số và phần thập phân có hai chữ số khác nhau. Số đó là $6{,}42$. Chữ số hàng phần trăm là bao nhiêu?
  \shortans{2}
  \loigiai{
    Chữ số hàng phần trăm là 2.
  }
\end{ex}

% Câu 20
\begin{ex}
  Làm tròn số $9{,}95$ đến hàng phần mười ta được số nào? (Nhân kết quả với 10).
  \shortans{100}
  \loigiai{
    Làm tròn thành $10{,}0$. Nhân 10 được 100.
  }
\end{ex}

% Câu 21
\begin{ex}
  Có bao nhiêu số thập phân có một chữ số ở phần thập phân nằm giữa $2$ và $3$?
  \shortans{9}
  \loigiai{
    Các số từ $2{,}1$ đến $2{,}9$. Có 9 số.
  }
\end{ex}

% Câu 22
\begin{ex}
  Viết số $5$ thành số thập phân có hai chữ số ở phần thập phân là $5{,}00$. Tổng các chữ số của số này là bao nhiêu?
  \shortans{5}
  \loigiai{
    $5 + 0 + 0 = 5$.
  }
\end{ex}

% Câu 23
\begin{ex}
  Đổi $1\text{ giờ } 12\text{ phút}$ thành giờ dưới dạng số thập phân là $1{,}2\text{ giờ}$. Chữ số hàng phần mười là bao nhiêu?
  \shortans{2}
  \loigiai{
    $12\text{ phút} = 0{,}2\text{ giờ}$. Số đó là $1{,}2\text{ giờ}$. Chữ số là 2.
  }
\end{ex}

% Câu 24
\begin{ex}
  Đổi $3\text{ km } 50\text{ m}$ thành ki-lô-mét dưới dạng số thập phân là $3{,}05\text{ km}$. Chữ số hàng phần mười là bao nhiêu?
  \shortans{0}
  \loigiai{
    Chữ số hàng phần mười là 0.
  }
\end{ex}

% Câu 25
\begin{ex}
  Tìm số tự nhiên $k$ biết $k$ là số lẻ và $14{,}8 < k < 18{,}2$. Tổng các giá trị của $k$ là bao nhiêu?
  \shortans{32}
  \loigiai{
    $k \in \{15, 17\}$. Tổng là $15 + 17 = 32$.
  }
\end{ex}
